\documentclass{amsart}
% For dealing with references we use the comment environment
\usepackage{verbatim}
\newenvironment{reference}{\comment}{\endcomment}
%\newenvironment{reference}{}{}
\newenvironment{slogan}{\comment}{\endcomment}
\newenvironment{history}{\comment}{\endcomment}

% For commutative diagrams we use Xy-pic
\usepackage[all]{xy}

% We use 2cell for 2-commutative diagrams.
\xyoption{2cell}
\UseAllTwocells

% We use multicol for the list of chapters between chapters
\usepackage{multicol}

% This is generally recommended for better output
\usepackage{lmodern}
\usepackage[T1]{fontenc}

% For cross-file-references

% Package for hypertext links:
\usepackage{hyperref}

% For any local file, say "hello.tex" you want to link to please

% Theorem environments.
%
\theoremstyle{plain}
\newtheorem{theorem}[subsection]{Theorem}
\newtheorem{proposition}[subsection]{Proposition}
\newtheorem{lemma}[subsection]{Lemma}

\theoremstyle{definition}
\newtheorem{definition}[subsection]{Definition}
\newtheorem{example}[subsection]{Example}
\newtheorem{exercise}[subsection]{Exercise}
\newtheorem{situation}[subsection]{Situation}

\theoremstyle{remark}
\newtheorem{remark}[subsection]{Remark}
\newtheorem{remarks}[subsection]{Remarks}

\numberwithin{equation}{subsection}

% Macros
%
\def\lim{\mathop{\mathrm{lim}}\nolimits}
\def\colim{\mathop{\mathrm{colim}}\nolimits}
\def\Spec{\mathop{\mathrm{Spec}}}
\def\Hom{\mathop{\mathrm{Hom}}\nolimits}
\def\Ext{\mathop{\mathrm{Ext}}\nolimits}
\def\SheafHom{\mathop{\mathcal{H}\!\mathit{om}}\nolimits}
\def\SheafExt{\mathop{\mathcal{E}\!\mathit{xt}}\nolimits}
\def\Sch{\mathit{Sch}}
\def\Mor{\mathop{\mathrm{Mor}}\nolimits}
\def\Ob{\mathop{\mathrm{Ob}}\nolimits}
\def\Sh{\mathop{\mathit{Sh}}\nolimits}
\def\NL{\mathop{N\!L}\nolimits}
\def\CH{\mathop{\mathrm{CH}}\nolimits}
\def\proetale{{pro\text{-}\acute{e}tale}}
\def\etale{{\acute{e}tale}}
\def\QCoh{\mathit{QCoh}}
\def\Ker{\mathop{\mathrm{Ker}}}
\def\Im{\mathop{\mathrm{Im}}}
\def\Coker{\mathop{\mathrm{Coker}}}
\def\Coim{\mathop{\mathrm{Coim}}}

% Boxtimes
%
\DeclareMathSymbol{\boxtimes}{\mathbin}{AMSa}{"02}

%
% Macros for moduli stacks/spaces
%
\def\QCohstack{\mathcal{QC}\!\mathit{oh}}
\def\Cohstack{\mathcal{C}\!\mathit{oh}}
\def\Spacesstack{\mathcal{S}\!\mathit{paces}}
\def\Quotfunctor{\mathrm{Quot}}
\def\Hilbfunctor{\mathrm{Hilb}}
\def\Curvesstack{\mathcal{C}\!\mathit{urves}}
\def\Polarizedstack{\mathcal{P}\!\mathit{olarized}}
\def\Complexesstack{\mathcal{C}\!\mathit{omplexes}}
% \Pic is the operator that assigns to X its picard group, usage \Pic(X)
% \Picardstack_{X/B} denotes the Picard stack of X over B
% \Picardfunctor_{X/B} denotes the Picard functor of X over B
\def\Pic{\mathop{\mathrm{Pic}}\nolimits}
\def\Picardstack{\mathcal{P}\!\mathit{ic}}
\def\Picardfunctor{\mathrm{Pic}}
\def\Deformationcategory{\mathcal{D}\!\mathit{ef}}

\setlength{\emergencystretch}{3em}
\begin{document}
\title[Stacks Project, Tag 0H9K]{406 --- Pro-isomorphisms preserve inverse limits and their first derived functors}
\maketitle
\noindent Copyright (C) 2005--2025 Johan de Jong. Stacks Project authors. Source: \href{https://stacks.math.columbia.edu/tag/0H9K}{Tag 0H9K}.

\noindent Editorial difficulty note (not author text). Difficulty H2: The proof constructs chain maps between product presentations using summed transition operators, rather than merely passing to limits. It then gives explicit coboundaries proving independence of a representative and invariance of the first derived limit under cofinal reindexing..

\noindent Editorial provenance note (not author text). History: excerpt from revision \texttt{a0444\allowbreak 6e57e\allowbreak c1fbc\allowbreak 252a8\allowbreak 71afc\allowbreak ec775\allowbreak 2fb28\allowbreak 07b14}, more-algebra.tex, lines 23072--23146. Reference presentation was adapted; original-author context and core support blocks were retained or added. The mathematical wording is unchanged. Licensed under GNU FDL 1.2 or later; no invariant sections or cover texts. See ../licenses/COPYING. Context blocks reproduce author definitions and supporting results; they are not additional counted problems.

\begin{lemma}
\label{lemma-pro-isomorphism-iso-Rlim}
Let $(A_n)$ and $(B_n)$ be inverse systems of abelian groups.
A morphism of pro-systems $\varphi : (A_n) \to (B_n)$ determines
maps $\lim A_n \to \lim B_n$ and $R^1\lim A_n \to R^1\lim B_n$.
These maps are isomorphisms if $\varphi$ is a pro-isomorphism.
\end{lemma}

\begin{proof}
Please see
Categories, Example \href{https://stacks.math.columbia.edu/tag/0G2W}{example pro morphism inverse systems (Tag 0G2W)}
for a discussion of morphisms of pro-systems.
The map $\varphi$ is given by $1 \leq m_1 < m_2 < m_3 < \ldots$
and maps $\varphi_n : A_{m_n} \to B_n$ compatible with transition
maps. Set $\varphi', \varphi'' : \prod A_n \to \prod B_n$
equal to
$\varphi'((a_n)) = (\varphi_n(a_{m_n}))$ and
$$
\varphi''((a_n)) =
(\varphi_n(a_{m_n} + f(a_{m_n + 1}) + \ldots + f(a_{m_{n + 1} - 1})))
$$
where each occurence of $f$ denotes a suitable transition map of
the inverse system $(A_n)$. Then the diagram
$$
\xymatrix{
\prod A_n \ar[r]_\delta \ar[d]^{\varphi'} &
\prod A_n \ar[d]^{\varphi''} \\
\prod B_n \ar[r]^\delta & \prod B_n
}
$$
where the horizontal arrows are as in Lemma \href{https://stacks.math.columbia.edu/tag/07KW}{lemma compute Rlim (Tag 07KW)}
is commutative. In this way we obtain the desired maps. The construction
is functorial in the sense that if we're given an inverse system $(C_n)$
and $1 \leq m'_1 < m'_2 < m'_3 < \ldots$ and maps $\psi_n : B_{m'_n} \to C_n$
compatible with transitition maps, then
$(\psi \circ \varphi)' = \psi' \circ \varphi'$ and
$(\psi \circ \varphi)'' = \psi'' \circ \varphi''$
where the composition $\psi \circ \varphi$ refers to the integers
$1 \leq m_{m'_1} < m_{m'_2} < \ldots$ and the maps
$\psi_n \circ \varphi_{m'_n} : A_{m_{m'_n}} \to C_n$. We will show
that if $B_n = A_n$ and $\varphi_n$ is the transition map, then
the resulting maps are the identity maps. This will both prove that
the construction is independent of the choice of the representative
$(m_n, \varphi_n)$ of $\varphi$ and the final statement of the lemma.

\medskip\noindent
Thus we let $B_n = A_n$ and $\varphi_n$ be the transition map.
Let $(a_n) \in \lim A_n$ be an element. Then $\varphi'((a_n))$
is the element which has in degree $n$ the image of $a_{m_n}$
which is equal to $a_n$. This proves the statement for $\lim A_n$.
Let $\xi \in R^1\lim A_n$ be the class of the element
$(a_n)$ in $\prod A_n$. Consider the element $(b_n)$ with
$$
b_i = a_i + f(a_{i + 1}) + \ldots + f(a_{m_n - 1})
$$
if $m_{n - 1} < i < m_n$ and $0$ if $i = m_n$ for some $n$.
Then $(a'_n) = (a_n) - \delta((b_n))$ defines the same class
in $R^1\lim A_n$ and a computation shows that $a'_i = 0$
unless $i = m_n$ for some $n$. Thus we may and do assume $a_i = 0$
unless $i = m_n$ for some $n$. Note that
$$
(0, \ldots, -f(a_{m_j}), 0, \ldots, 0, a_{m_j}, 0, \ldots)
$$
with nonzero entries in spots $j$ and $m_j$, is the image under $\delta$ of
$$
c_j = (0, \ldots, 0, f(a_{m_j}), f(a_{m_j}), \ldots, a_{m_j}, 0, \ldots)
$$
with nonzero entries in spots $j + 1, \ldots, m_j$. The sum
$c = \sum c_j$ makes sense in $\prod A_n$. Recalling that
$\varphi''((a_n)) = (a_{m_1}, a_{m_2}, \ldots)$ we see that
$$
(a_n) - \varphi''((a_n)) = \delta(c)
$$
and the proof is complete.
\end{proof}
\end{document}
